% GENERATED FILE. Edit canonical lessons, then run npm run generate:books.
% canonical-source-hash: fnv1a64:f5a1cb6592ad40f6
% canonical-lessons: UR-C78-vaqt, UR-C78-dur, UR-C78-qarib, UR-C78-andar, UR-C78-bahar

\chapter{Time, Far, Near, Inside, Outside}
\label{ch:ur-time-far-near-inside-outside}
\hlchaptermodality{78}

\begin{chapteropening}
\textbf{By the end of this chapter:} \emph{I can say time, far, near, inside and outside.}

Complete the last lesson of chapter 78: I can say time, far, near, inside and outside.
\end{chapteropening}

\section[\texorpdfstring{\ur{وقت} (vaqt)}{وقت (vaqt)}]{\ur{وقت} (vaqt) — time}

\label{lesson:UR-C78-vaqt}



\begin{quote}
Before the new one: say the Urdu for a delay; late, then the Urdu for before.
\end{quote}

\subsection*{You'll want to know: \ur{وقت}}
\textbf{\ur{وقت}} — \emph{vaqt} — \textquotedblleft{}time\textquotedblright{}.

\begin{cousinweb}
\textbf{\ur{وقت}} (\emph{vaqt}) is a loan from Arabic \emph{waqt}, \textquotedblleft{}time, an appointed time\textquotedblright{}, which reached Urdu through Persian. Its Arabic root is \emph{w-q-t} (Lane).
\end{cousinweb}

\begin{grammarlens}[title={What you've built}]
One more word for this chapter.
\end{grammarlens}

\subsection*{Your turn}
\begin{itemize}
\raggedright
  \item \emph{Say it:} \emph{vaqt}
  \item \emph{Say it:} \emph{vaqt}, once more
  \item \emph{Say it:} say \emph{vaqt}
  \item \emph{Recall:} say the Urdu for a delay; late, then the Urdu for before, then say \emph{vaqt} again
\end{itemize}

\begin{tcolorbox}[breakable,colback=teal!4,colframe=teal!35!black,title={Before you move on}]
What does \ur{وقت} mean? (\textquotedblleft{}Time\textquotedblright{}.) Say it once more.
\end{tcolorbox}
\section[\texorpdfstring{\ur{دور} (dūr)}{دور (dūr)}]{\ur{دور} (dūr) — far}

\label{lesson:UR-C78-dur}



\begin{quote}
Before the new one: say the Urdu for before, then the Urdu for time.
\end{quote}

\subsection*{You'll want to know: \ur{دور}}
\textbf{\ur{دور}} — \emph{dūr} — \textquotedblleft{}far\textquotedblright{}.

\begin{cousinweb}
\textbf{\ur{دور}} (\emph{dūr}) is a tadbhava, a word inherited from Sanskrit and reshaped by sound change. It comes from \emph{dūra}, \textquotedblleft{}distant\textquotedblright{} (CDIAL 6495).
\end{cousinweb}

\begin{grammarlens}[title={What you've built}]
One more word for this chapter.
\end{grammarlens}

\subsection*{Your turn}
\begin{itemize}
\raggedright
  \item \emph{Say it:} \emph{dūr}
  \item \emph{Say it:} \emph{dūr}, once more
  \item \emph{Say it:} say \emph{dūr}
  \item \emph{Recall:} say the Urdu for before, then the Urdu for time, then say \emph{dūr} again
\end{itemize}

\begin{tcolorbox}[breakable,colback=teal!4,colframe=teal!35!black,title={Before you move on}]
What does \ur{دور} mean? (\textquotedblleft{}Far\textquotedblright{}.) Say it once more.
\end{tcolorbox}
\section[\texorpdfstring{\ur{قریب} (qarīb)}{قریب (qarīb)}]{\ur{قریب} (qarīb) — near}

\label{lesson:UR-C78-qarib}



\begin{quote}
Before the new one: say the Urdu for time, then the Urdu for far.
\end{quote}

\subsection*{You'll want to know: \ur{قریب}}
\textbf{\ur{قریب}} — \emph{qarīb} — \textquotedblleft{}near\textquotedblright{}.

\begin{cousinweb}
\textbf{\ur{قریب}} (\emph{qarīb}) is a loan from Arabic \emph{qarīb}, \textquotedblleft{}near\textquotedblright{}, which reached Urdu through Persian. Its Arabic root is \emph{q-r-b} (Lane).
\end{cousinweb}

\begin{grammarlens}[title={What you've built}]
One more word for this chapter.
\end{grammarlens}

\subsection*{Your turn}
\begin{itemize}
\raggedright
  \item \emph{Say it:} \emph{qarīb}
  \item \emph{Say it:} \emph{qarīb}, once more
  \item \emph{Say it:} say \emph{qarīb}
  \item \emph{Recall:} say the Urdu for time, then the Urdu for far, then say \emph{qarīb} again
\end{itemize}

\begin{tcolorbox}[breakable,colback=teal!4,colframe=teal!35!black,title={Before you move on}]
What does \ur{قریب} mean? (\textquotedblleft{}Near\textquotedblright{}.) Say it once more.
\end{tcolorbox}
\section[\texorpdfstring{\ur{اندر} (andar)}{اندر (andar)}]{\ur{اندر} (andar) — inside}

\label{lesson:UR-C78-andar}



\begin{quote}
Before the new one: say the Urdu for far, then the Urdu for near.
\end{quote}

\subsection*{You'll want to know: \ur{اندر}}
\textbf{\ur{اندر}} — \emph{andar} — \textquotedblleft{}inside\textquotedblright{}.

\begin{cousinweb}
\textbf{\ur{اندر}} (\emph{andar}) is a loan, usually traced to Persian \emph{andar}, \textquotedblleft{}inside\textquotedblright{}.
\end{cousinweb}

\begin{grammarlens}[title={What you've built}]
One more word for this chapter.
\end{grammarlens}

\subsection*{Your turn}
\begin{itemize}
\raggedright
  \item \emph{Say it:} \emph{andar}
  \item \emph{Say it:} \emph{andar}, once more
  \item \emph{Say it:} say \emph{andar}
  \item \emph{Recall:} say the Urdu for far, then the Urdu for near, then say \emph{andar} again
\end{itemize}

\begin{tcolorbox}[breakable,colback=teal!4,colframe=teal!35!black,title={Before you move on}]
What does \ur{اندر} mean? (\textquotedblleft{}Inside\textquotedblright{}.) Say it once more.
\end{tcolorbox}
\section[\texorpdfstring{\ur{باہر} (bāhar)}{باہر (bāhar)}]{\ur{باہر} (bāhar) — outside}

\label{lesson:UR-C78-bahar}



\begin{quote}
Before the new one: say the Urdu for near, then the Urdu for inside.
\end{quote}

\subsection*{You'll want to know: \ur{باہر}}
\textbf{\ur{باہر}} — \emph{bāhar} — \textquotedblleft{}outside\textquotedblright{}.

\begin{cousinweb}
\textbf{\ur{باہر}} (\emph{bāhar}) is a tadbhava, a word inherited and reshaped by sound change. It comes from \emph{*bāhira}, \textquotedblleft{}external\textquotedblright{}, a form CDIAL reconstructs from its descendants (CDIAL 9226).
\end{cousinweb}

\begin{grammarlens}[title={What you've built}]
One more word for this chapter.
\end{grammarlens}

\subsection*{Your turn}
\begin{itemize}
\raggedright
  \item \emph{Say it:} \emph{bāhar}
  \item \emph{Say it:} \emph{bāhar}, once more
  \item \emph{Say it:} say \emph{bāhar}
  \item \emph{Recall:} say the Urdu for near, then the Urdu for inside, then say \emph{bāhar} again
\end{itemize}

\begin{tcolorbox}[breakable,colback=teal!4,colframe=teal!35!black,title={Before you move on}]
What does \ur{باہر} mean? (\textquotedblleft{}Outside\textquotedblright{}.) Say it once more.
\end{tcolorbox}
